\section{Models Development}\label{sec:models_development}

The model development procedure were divided into two parts: baseline model development and then advanced model. 

\subsection{Baseline Model}\label{sec:baseline_model}

The baseline model is built using only one predictor to assess the overall performance, using only one of the most essential variables in predicting fuel, namely Speed. 

At first, Speed over ground was selected to model it with fuel consumption, then the engine rpm. The study revealed that using engine rpm is a stronger predictor than using Speed, as shown in the Results section. (We could filter out Speed based on the technical documentation of the engine to only consider valid speed limits.)

Next the outliers in the response variable have been dropped by integration to the preprocessing pipeline which reduced the overall dataset to xx observations so xx percentage of the overall dataset. This has proven to be beneficial for the overall model accuracy as discussed in Discussion section. 

Afterward, different standardization techniques such as standard scaling, normalization, and min-max scaling (perhaps more info about those techniques?) have were employed. However, they didn't improve the accuracy, and therefore, they were not considered. 

Then usage of the fuel consumption of the main engine instead of total fuel consumption which is combination of main engine, boiler and auxiliary consumption have been compared. The insights discussed in the Discussion section explained the need to use only the main engine fuel instead of total fuel, and therefore, this was selected. 

Then, different estimators, linear and non-linear categories were used fitted, and compared against each other using five-fold cross-validation (more about cross-validation?), and the following performance metrics were used to score the models based on the average for the folds and standard deviations $R^2$ and 
$\text{RMSE} = \sqrt{\frac{1}{n} \sum_{i=1}^{n} (y_i - \hat{y}_i)^2}$

K-fold cross-validation is a widely adopted resampling technique in machine learning that ensures model generalization to unseen data. It uses different subsets of the data for training and testing and then averages the final results across the folds. It's a reliable performance estimate that balances the bias/variance tradeoff and ensures less variability and more use of small datasets. 
The estimators used with baseline configuration are depicted in the sklearn pipeline shown in **Fig. 9**. 

\begin{wrapfigure}{l}{0.6\textwidth}
    \centering
    \includegraphics[width=1\linewidth]{baseline_pipeline}
    \caption{Baseline Pipeline}
    \label{fig:baseline_pipeline}
\end{wrapfigure}

(original as is from prediction using black and white box paper)
The black-box models, as a type of regression model, can be evaluated using mean square error \ac{mse} (MSE), \ac{rmse} root-mean-square error (RMSE), \ac{mae} mean absolute error (MAE), MAPE, and $R^2$ \ac{r2}.

The performance metrics used were the $R^2$, namely the Coefficient of Determination, which measures the variance in the dependent variable explained by the independent variables. 

\[
R^2 = 1 - \frac{\sum_{i=1}^{n} (y_i - \hat{y}_i)^2}{\sum_{i=1}^{n} (y_i - \bar{y})^2}
\]

\begin{itemize}
    \item \( y_i \): The actual observed values of the dependent variable in the dataset.
    \item \( \hat{y}_i \): The predicted values of the dependent variable based on the model.
    \item \( \bar{y} \): The mean (average) of all actual values \( y_i \); represents the central tendency of the observed data.
    \item \( n \): The number of observations or data points in the dataset.
    \item \( \sum_{i=1}^{n} (y_i - \hat{y}_i)^2 \): The sum of squared residuals (errors) between actual and predicted values; indicates the unexplained variance by the model.
    \item \( \sum_{i=1}^{n} (y_i - \bar{y})^2 \): The total sum of squares, which represents the total variance in the observed data.
\end{itemize}

Therefore, adjusted $R^2$ comes into play for the advanced models, a modified version of $ R^2$ that adjusts for the number of predictors in the model. It penalizes the addition of irrelevant variables that do not improve the model, making it more suitable for comparing models with different numbers of predictors.

\[
\text{Adjusted } R^2 = 1 - \left(1 - R^2\right) \frac{n - 1}{n - k - 1}
\]

\begin{itemize}
    \item \( R^2 \): The coefficient of determination, representing the proportion of variance in the dependent variable that is explained by the independent variables.
    \item \( n \): The total number of observations in the dataset.
    \item \( k \): The number of independent variables (predictors) in the model.
\end{itemize}

\ac{rmse} measures the average magnitude of the errors between the predicted and observed values. It gives an idea of how concentrated the data is around the line of best fit.

\[
\text{RMSE} = \sqrt{\frac{1}{n} \sum_{i=1}^{n} (y_i - \hat{y}_i)^2}
\]

\begin{itemize}
    \item \( n \): The total number of observations in the dataset.
    \item \( y_i \): The actual observed value for the \(i\)-th observation.
    \item \( \hat{y}_i \): The predicted value for the \(i\)-th observation from the model.
    \item \( (y_i - \hat{y}_i)^2 \): The squared difference (error) between the actual and predicted values for each observation.
\end{itemize}

Next different seeds have been search to accomodate for the error due to the random shuffle procedure, this has been shown in Results section. 

\subsection{Advanced Model}\label{sec:advanced_model}

Next, an advanced multivariate model was developed and compared against a baseline model in which all the selected predictors were utilized. This has revealed a tangible improvement in overall performance. 

Next, the most promising model was selected, and the grid searched for the best hyperparameters and most predictable features. The combination of the best parameters and features subset was then used to train the final model.

- All features available from voyage reports vs baseline 
- With ocean features vs without 
- with climate features vs without 
- Waves data from cmems vs era5 

also, the data and parameters used 

The trained model is validated using a k-fold cross-validation procedure. This ensures accuracy and reliability.